\scanpage{73}
\mspage{12}
dont le deuxième est donné par
\[
\text{(41)}\qquad
\boxed{\rho\longmapsto\rho_\infty^{-1},\quad\sigma_\infty\longmapsto\rho_1},
\qquad
\underset{\substack{\Vert\\\sigma_\infty\rho}}{\varepsilon_0}\longmapsto\rho_0
\]
\margin{Attention à l'exposant !}
et le \uline{premier par}
\[
\begin{aligned}
\text{(42)}\qquad
&\boxed{\tau_\infty\longmapsto\tau_\infty},\qquad
\underset{\substack{\Vert\\\rho(\tau_\infty)}}{\tau_0}
\longmapsto
\underbrace{\overbrace{\rho_\infty^{-1}\tau_\infty\rho_\infty}^{\tau_0\tau_1\tau_\infty\tau_1\tau_0}}_{(\rho_\infty^{-1}\rho_1\rho_\infty^{-1}\rho_1)\tau_\infty},\\
&
\underset{\substack{\Vert\\\rho^2(\tau_\infty)}}{\tau_1}
\longmapsto
\overbrace{\rho_\infty^{-2}\tau_\infty\rho_\infty^2
=\rho_\infty\tau_\infty\rho_\infty^{-1}}^{\substack{(=(\rho_\infty\rho_1\rho_\infty\rho_1)\tau_\infty)\\
\overbrace{\tau_1\tau_0\tau_\infty\tau_0\tau_1}}}
\end{aligned}
\]
\margin{(car $\rho_\infty^3=1$ dans $\Pi_{03}^{\tau}/\langle\ldots\rangle$)}
\margin{\uline{NB} On utilise $\rho_1^2=\rho_\infty^3=1$ et
$\tau_\infty(\rho_\infty)=\rho_0\rho_1=\rho_1^{-1}\rho_\infty^{-1}\rho_1=\rho_1\rho_\infty^{-1}\rho_1$.}

où
\[
\rho,\sigma_i,\varepsilon_i\quad(i\in\{0,1,\infty\})
\quad\text{dans }\Pi_1(U_{0,3},\unsure{\mathcal G_3|\mathfrak S_3};P_+)=\Pi_{0,3}^{D_3}
\]
sont définis de façon \unsure{standart} (on y reviendra dans un cas plus général un peu plus bas), et
\[
\tau_i\ (i\in\{0,1,\infty\})\in
\underbrace{\Pi_1(U_{0,3},\{1,\tau\},P_+)}_{\Pi_{0,3}^{\tau}}
\subset
\underbrace{\Pi_1(U_{0,3},\unsure{\mathcal G_3|\mathfrak S_3}\times\{1,\tau\};P_+)}_{\Pi_{0,3}^{D_3,\tau}}
\]
sont définis dans (24) (ce sont les générateurs de $\Pi_{0,3}^{\tau}\simeq\unsure{\mathfrak G_2}$ qui figurent dans (2)). Je vais revenir sur ces isomorphismes (40) dans un contexte un peu plus général, ci-dessous.

Enfin, dans le développement précédent, remplaçant $U_{0,3}$ et l'action de $\unsure{\mathcal G_3|\mathfrak S_3}$ resp. $\unsure{\mathcal G_3|\mathfrak S_3}\times\{1,\tau\}$ dessus, par \deleted{celle de} le demi plan de Poincaré $\Sigma^+$ et l'action de $\unsure{\operatorname{Sl}|\operatorname{Gl}}(2,\mathbb Z)/\pm1$ resp. $\operatorname{Gl}(2,\mathbb Z)/\pm1$ dessus (par l'action homographique resp. homographique et ``antihomographique''), on trouve enfin des isomorphismes remarquables
\diagram{\begin{tikzcd}[ampersand replacement=\&,column sep=large]
\text{(43)}\quad \operatorname{Gl}(2,\mathbb Z)/\pm1 \arrow[r,"\sim"]
\& \Pi_{0,3}^{D_3,\tau}\quad(
\arrow[r,dash,"\sim"]
\&\Pi_{0,3}^{\tau}/\langle\rho_1^2,\rho_\infty^3\rangle)
\\
\text{(44)}\quad \operatorname{Sl}(2,\mathbb Z)/\pm1 \arrow[r,"\sim"]
\& \Pi_{0,3}^{D_3}\quad(
\arrow[r,"\sim"]
\&\Pi_{0,3}/\langle\rho_1^2,\rho_\infty^3\rangle\unsure{)}
\end{tikzcd}}
\margin{72}
